Un’azienda riceve ogni giorno decine di ticket via email o modulo web: richieste su fatture e pagamenti (Amministrazione), segnalazioni di malfunzionamenti (Tecnico), domande su offerte e ordini (Commerciale). Lo smistamento manuale rallenta i tempi di risposta e crea frustrazione.
\\\\L’obiettivo del project work è realizzare un prototipo essenziale e riproducibile che, a partire da un piccolo dataset sintetico di ticket brevi (oggetto + descrizione), classifichi automaticamente ogni ticket nella categoria più pertinente e suggerisca una priorità operativa (bassa/media/alta).
\\\\Si richiede l’uso di tecniche di ML di base mantenendo il focus su semplicità del codice e chiarezza del flusso: input testuale, output con categoria e priorità, e una misura elementare di qualità (ad esempio accuracy o F1) per verificare che il modello sia ragionevolmente affidabile. Tutti i dati devono essere generati ad hoc, senza informazioni personali, privilegiando un design minimale che permetta di eseguire e comprendere il prototipo in pochi minuti.